\begin{table}[htbp]
\centering
\small
\caption{Sample eligibility. The common comparison set is defined by M4, the
most demanding model. Exclusions are not mutually exclusive; the union is 229
cells. Spatial statistics use 2,535 observations after removing one Queen
contiguity island.}
\label{tab:eligibility}
\begin{threeparttable}
\begin{tabularx}{\textwidth}{@{}>{\raggedright\arraybackslash}p{6.1cm} r >{\raggedright\arraybackslash}X@{}}
\toprule
Category & Cells & Note \\
\midrule
Master grid & 2,765 & Five mainland cities; 500 m lattice \\
M1 / M2 eligible & 2,764 & Excludes one low-LST-support cell (\texttt{JKT\_R037\_C062}) \\
M3 eligible & 2,758 & Also excludes six cells without population support \\
M4 eligible (common set) & 2,536 & Also excludes age-unsupported and route-missing cells \\
Excluded from common set & 229 & 8.28\% of master grid \\
\quad Route missing only & 222 & No path to a candidate park in baseline graph \\
\quad Route + population + age missing & 4 & Already outside common set because of route \\
\quad Population/age missing only & 2 & \texttt{JKT\_R003\_C031}, \texttt{JKT\_R057\_C063} \\
\quad Low LST support & 1 & \texttt{JKT\_R037\_C062}; 27.759\% valid LST area \\
Spatial observations after island removal & 2,535 & \texttt{JKT\_R013\_C019} is a Queen island \\
\bottomrule
\end{tabularx}
\begin{tablenotes}[flushleft]\footnotesize
\item The two population-only exclusions explain the difference between the 2,538-cell benchmark and the 2,536-cell common set. They are slivers (0.000266 and 0.002643 km$^2$) with no source support; their original exported zeros are preserved as raw values, not treated as estimated population.
\end{tablenotes}
\end{threeparttable}
\end{table}
